\Problem{Números bicolor}

\Statement

Definimos como número \emph{bicolor} un natural $n$ con solo dos dígitos diferentes que están repetidos en dos bloques (o "colores"). Más formalmente, la secuencia de dígitos de $n$ es $d_1d_2\ldots d_ke_1e_2\ldots e_l$, donde $d$ y $e$ son los dos dígitos y, $d\neq e$, $k > 0$ y $l > 0$.

\medskip

Por ejemplo, son números bicolor: 7722, 44111, 666699, 277 y 45. \\
Y los siguiente \textbf{no} son bicolor: 121, 113311, 7878, 1234, 7, 55, 99910.

\medskip

Implementa una \textbf{función} \verb|is_bicolor| que recibe un natural y determina si es bicolor. La función recibe un número $n > 0$ y devuelve \texttt{true} si es bicolor y \texttt{false} en caso contrario.

\medskip

La cabecera de la función debe ser, exactamente:

\begin{verbatim}
/**
* @pre n >= 0
* @post retorna cierto si n es bicolor, falso en caso contrario
*/
bool is_bicolor(int n);
\end{verbatim}

\Observation

Solo hay que enviar la función solicitada; el programa principal será ignorado.

